%% 04-methodology.tex --- included by ../sfc.tex
\section{Experimental Methodology}
\label{sec:methodology}

\subsection{Research Questions}
\label{sec:rqs}

The evaluation answers four questions, each with its own controls, and
Section~\ref{sec:results} is organised around them.

\begin{itemize}[leftmargin=1.8em,itemsep=1pt,topsep=2pt]
\item[\textbf{RQ1}] \textbf{End-to-end cost.} Under each compiler's own
  recommended protocol, how does \PORTER{}'s EPR consumption and compile
  time compare with \TeleSABRE{}'s and \dSABRE{}'s?
\item[\textbf{RQ2}] \textbf{Routing quality with the layout controlled.}
  Given the same, possibly externally supplied, initial layout, how does
  \PORTER{}'s routing loop compare with \dSABRE{}'s --- first at an equal
  route count, then when each is allowed its own refinement schedule?
\item[\textbf{RQ3}] \textbf{Mechanism attribution.} What does each of
  \PORTER{}'s mechanisms --- capacity relief, the packet primitive, the
  \TeleGate{} family, the initial layout and horizon selection ---
  contribute on its own?
\item[\textbf{RQ4}] \textbf{Robustness.} How do the answers change with
  circuit width, core count and core-graph shape?
\end{itemize}

Each results subsection opens with the question it answers, what is held
fixed, and the metric; the interpretation of the answers is collected in
Section~\ref{sec:discussion}, and the assumptions and limitations they
rest on in \appThreats{}.


\subsection{Terminology}
\label{sec:terms}

Budget comparisons depend on using these terms consistently, so we fix
them here and use no synonyms.

\begin{itemize}[leftmargin=1.2em,itemsep=0pt,topsep=2pt]
\item A \textbf{route} (equivalently a \emph{routing pass}) is one
  invocation of the routing loop of Section~\ref{sec:sfc} on one circuit
  from one layout, producing one operation stream.
\item A \textbf{candidate layout} is one initial mapping a protocol
  supplies to the router. A protocol's cost is its layout count times its
  routes per layout.
\item A \textbf{protocol} is a layout source, a route budget and a
  selection rule; Table~\ref{tab:budgets} names the ones used in the
  main results. 
\item A \textbf{candidate horizon} is one value of $L_\mathrm{tp}$ offered
  to horizon selection; selecting among them costs one route each.
\item \textbf{Complete compile cost} is the wall-clock time of an entire
  protocol --- layout generation plus every route over every candidate ---
  not of the winning route alone.
\end{itemize}

For the machine and the router we use: \emph{core}, \emph{communication
port}, and \emph{reservoir} (a core holding at least $f_\mathrm{min}$ free
positions); a front CX is \emph{local}, \emph{adjacent-remote} or
\emph{non-adjacent-remote}; a plan is \emph{normal} or \emph{blocked}, and
\emph{ready}, \emph{preparable} or neither. A run that does not complete
does so in one of four ways, distinguished throughout: a capacity
\emph{deadlock} leaves no legal plan and no relief path, and raises a
diagnostic; a \emph{livelock} makes local progress indefinitely while
retired gates and EPR count stand still, and is stopped by the
no-progress guard; a \emph{timeout} exceeds the wall-clock cap; and an
\emph{infeasible input} is one the machine has no spare capacity for,
refused before routing begins (Section~\ref{sec:sfc-progress}).

\subsection{Setup}
\label{sec:setup}
\label{sec:perf}

\textbf{Implementation.} \PORTER{} is Python, released with the harness.
It takes no seed: heuristic ties resolve through a fixed chain of
deterministic keys and any residual tie through a constant-initialised
generator, so one route on a fixed circuit and hardware emits a
byte-identical stream across \chg{runs, horizon selection included: a
speculative route is bounded by the DAG layers its lookahead walk may peel,
at most seven times the count of the route the configured horizon produced,
so the bound is a property of the route rather than of the machine. The
horizon each run selected is recorded with its result. Every reported \PORTER{} route is certified by}
a replay verifier written separately from the router, which re-derives
gate coverage, dependency order and physical legality from the emitted
stream alone and is itself tested by fault injection
(\appVerification{}).

\textbf{Architectures.} We use the two grid-of-grids topologies
\dSABRE{}~\cite{lidsabre} evaluates on --- architecture~A, a $2{\times}2$
array of $3{\times}3$ cores (36 physical qubits), and architecture~H, a
$2{\times}3$ array of $4{\times}4$ cores (96 physical) --- adding its
$5{\times}5$-core scalability series in Section~\ref{sec:scaleup} and two
networks of IBM 27-qubit heavy-hex tiles in Section~\ref{sec:heavyhex}.
All are read from the baselines' own device files, so every compiler sees
identical physical qubit ids.

\textbf{Benchmarks.} We use the same
MQT-Bench~\cite{quetschlich2023mqtbench} files \dSABRE{}'s evaluation
uses, confirmed by exact CX-count match: a 25-qubit suite (AE, GHZ,
Graphstate, QFT, QNN, Random) on architecture~A, and a 64-qubit suite
adding QPEexact, QAOA and Multiplier on architecture~H. We add an
80-qubit suite of the same nine families; its AE and QNN instances are
not produced by the stock generator and were reconstructed to match its
construction.

\label{sec:baselines}
\textbf{Baselines.} We compare against
\TeleSABRE{}~\cite{russo2025telesabre} and \dSABRE{}~\cite{lidsabre}, the
latter in its default configuration --- safe mode, which makes capacity a
legality condition with a per-core reserve and recovers by relaying slack
along a BFS path (Section~\ref{sec:related}).

Three decisions avoid disadvantaging \dSABRE{}. Its recovery budgets
derive from the architecture as its own drivers do, since a binding
iteration budget would cut its heuristic search short and hand the
remainder to its guaranteed drain at a worse EPR count. We do not charge
it for pre-routing layout repair, which lands in its EPR counter although
it adjusts the starting layout rather than routing anything; we net that
out for both compilers. And we take its layout rule exactly as it ships
(Section~\ref{sec:protocols}).

\TeleSABRE{}'s stock C binary crashes deterministically on
architecture~A, in its initial-layout pass rather than in routing. All
\TeleSABRE{} numbers use a minimally patched binary, released with ours,
that falls back to the least-loaded core; it moves only qubits the stock
binary leaves unplaced, and architecture-H results are unchanged. Every
compiler invocation is capped at $900$\,s wall-clock, and a call
exceeding it counts as a failure to complete.

\subsection{Initial-Layout Protocols}
\label{sec:protocols}

We evaluate under three initial-layout protocols, designed to separate
routing quality from initial-layout quality.

\textbf{Own-protocol comparison.} Each router is first evaluated under
its own established best-practice protocol. For the two baselines that
is best-of-three, matching the convention already used to compare
\TeleSABRE{} and \dSABRE{} to each other~\cite{lidsabre}:
\TeleSABRE{} runs its built-in Hungarian initial-layout heuristic at
config seeds 0--2, and \dSABRE{} runs Qiskit \texttt{SabreLayout} on
its corners-removed architecture graph at seeds 0--2, each layout
routed by \SABRE{}-style forward\,$\to$\,backward\,$\to$\,forward
passes; both report the best EPR count. \PORTER{} instead builds one
deterministic temporal locality-greedy layout
(Section~\ref{sec:sfc-layout}) and routes it at each of the three
candidate lookahead horizons of Section~\ref{sec:horizon}, keeping the
best --- \chg{four routing passes over one layout, against nine over three.}
It searches no layouts, so every number is reproducible from the circuit
file alone.

\dSABRE{}'s SabreLayout runs on a coupling map with the most remote
minimum-degree (corner) physical qubits of every core removed, the count
per core fixed by its own fill-adaptive rule. Removing them from the map
means no logical qubit is placed there, so they begin free and hold spare
capacity for arriving teleports and evictions; they are not communication
ports, which are the designated qubits at the link endpoints
(Section~\ref{sec:qcircuits}). We use the rule as it ships, which is why
our \dSABRE{} figures reproduce those it publishes for the architectures
it reports.

\textbf{Shared-layout comparison.} To separate layout quality from
routing quality, both \dSABRE{} and \PORTER{} are then run from an
identical externally supplied mapping --- \TeleSABRE{}'s own initial
layout, and separately \texttt{pytket-dqc}'s hypergraph partition
(\appPytket{}) --- so the only difference left is the routing loop.
\textbf{Cross-layout comparison.} Finally each compiler routes the
other's own layout, which closes the remaining cell of the design.

\subsection{Operating Protocols}
\label{sec:opprotocols}

Section~\ref{sec:sfc} defines one routing pass. We evaluate two protocols
built on it. \PORTER{} routes its own layout once at each of three
candidate lookahead horizons --- \chg{$L_\mathrm{tp}\in\{4,8,16,128\}$, one
fixed set applied to every circuit --- and keeps the lowest-EPR result; \PORTER{}+R refines backward}
from that selected route. Both are
deterministic and both keep the best result by EPR, so neither can return
a worse route than the pass it contains. They invoke the router
repeatedly rather than changing what a pass does, which is why they are
described here and not in Section~\ref{sec:sfc}.

\subsubsection{\PORTER{}: choosing the horizon per circuit}
\label{sec:horizon}

The teleport-lookahead horizon $L_\mathrm{tp}$ is the one parameter of
Section~\ref{sec:sfc-scoring} for which no constant is best across the
corpus. In an offline sweep over $\{4,8,\dots,1024\}$ the best value
spans the whole range --- $80$-qubit QNN wants $4$, the $25$-qubit suite
$16$, every $100$-, $200$- and $360$-qubit instance $128$ --- and the gap
between \PORTER{}'s configured $L_\mathrm{tp}{=}16$ and the per-instance
best is $17.2\%$ geometric-mean EPR.

Interaction density sets the direction. Let $\rho$ be the mean, over
logical qubits appearing in at least one CX, of the distinct partners
each shares a CX with anywhere in the circuit, divided by $n-1$: one
scalar summarising the whole circuit, not a per-qubit or horizon-limited
quantity. So $\rho\to1$ is all-to-all and a linear CX chain gives
$\rho\approx2/(n-1)$. It is measured on the transpiled benchmark
rather than the textbook algorithm, and the two can differ sharply:
optimisation level 3 drops the controlled-phase rotations whose angle
underflows its tolerance, which leaves QFT between $30$ and $37$ partners
per qubit at every width above $25$, so its $\rho$ falls from $0.92$ at
$25$ qubits to $0.10$ at $360$. A circuit family need not stay on one
side of the threshold as it scales. Near $\rho{=}1$ a long
horizon finds the moving qubit wanted in every core at once, the gain
estimate degenerates into noise, and short horizons win; where $\rho$ is
small the signal is coherent and long horizons win. It is not a usable
predictor alone, though. Its rank correlation with the per-instance best
horizon is $-0.765$ (Spearman), strong but not decisive: in the band
$0.42\le\rho\le0.53$ where most 64-qubit circuits sit, QAOA wants $8$ and
AE wants $128$, and a threshold rule on $\rho$ recovers only $8.0\%$ of
the available $17.2\%$. That one scalar cannot separate them is
unsurprising --- the horizon acts per moving qubit, while $\rho$ averages
over all of them --- and a per-qubit or horizon-limited density is
untried.

\PORTER{} therefore \chg{does not use $\rho$ at all. Every circuit is routed
at the same four horizons, $L_\mathrm{tp}\in\{4,8,16,128\}$, and the
lowest-EPR route is kept --- four routes over one layout. We prefer this to a
rule keyed on $\rho$, or on any other circuit statistic, because such a rule is
fitted on the same instances the evaluation reports and buys nothing: measured
over all forty-eight instances the fixed set matches a tuned two-set rule on the
main suite and is $0.2\%$ better pooled, and under backward refinement it is
identical on every one of the twenty-four main instances. What the rule was
buying was compile time, not EPR --- it kept the long horizon away from
near-all-to-all circuits, where it does not converge --- and the fixed set pays
that back at $1.6$ times the compile time for one pass and $1.3$ with
refinement. Three values are not enough: the best three-value set is $0.3\%$
worse than the rule. Nor is this particular quartet distinguished;
$\{4,16,32,128\}$, $\{4,16,48,128\}$ and $\{4,16,64,128\}$ all land within a
percent of it, and the last is slightly better pooled. We take
$\{4,8,16,128\}$ because it is the union of the two sets the retired rule chose
between, so it introduces no horizon value the compiler did not already use.
Two rules bound}
the cost: the configured horizon --- $L_\mathrm{tp}{=}16$, the default of
Section~\ref{sec:sfc-scoring}, which \chg{the set contains --- is
routed first, unguarded, and wins ties; each other candidate is bounded by
the DAG layers its lookahead walk may peel, at most seven times what that first
route peeled. Hence a}
conditional guarantee --- \emph{provided the configured-horizon route
completes}, selection never returns a higher EPR count than that route
\chg{alone.}

\chg{\textbf{The candidate set was fixed on the corpus it is evaluated on.}
The sweep and the four values were chosen by inspecting the $48$ instances,
with no held-out family, so evidence for the \emph{size} of the gain is
in-sample. What that risks is bounded, and less than a selection rule would
risk: the set is a portfolio rather than a prediction, so a poor choice costs
compile time and cannot worsen the selected route; the configured horizon is a
member and wins ties, so no instance can do worse than routing once; and
because nothing is conditioned on the circuit, there is no threshold to
generalise. Several other four-value sets perform within a percent, which is
the evidence that the result does not turn on these particular numbers. An}
out-of-sample validation remains to be done.

\subsubsection{\PORTER{}+R: refining backward from that route}
\label{sec:sfc-refine}

\dSABRE{} routes each candidate layout forward, backward and forward
again, so comparing route budgets rather than layout counts alone needs a
\PORTER{} protocol that spends the same effort. \PORTER{}+R is
\SABRE{}'s bidirectional layout refinement~\cite{li2019tackling} applied
to whole routes. Write $R(C,\ell)$ for routing circuit $C$ from layout
$\ell$, and $\ell'_i$ for the layout pass $i$ ends on. Pass~1 computes
$R(C,\ell_0)$ and ends on $\ell'_1$; pass~2 computes
$R(C^{\mathrm{rev}},\ell'_1)$ on the reversed circuit, \emph{discards its
route} and keeps only $\ell'_2$; pass~3 computes $R(C,\ell'_2)$.
Reversing $C$ reverses every dependency edge, so routing it from where
the forward pass ended yields a layout adapted to the \emph{start} of the
forward circuit. Odd pass counts above three repeat the cycle;
$P{=}5$ throughout, the best forward pass is kept, and nothing is random.

Pass~1 \emph{is} \PORTER{}: it runs the selection above, and the horizon
it chooses is held for the remaining $P-1$ passes --- re-selecting each
pass would multiply the count by the number of candidates to re-decide a
property refinement does not change. That is what makes \PORTER{}+R
contain \PORTER{}, so with the best forward pass kept it can never report
more EPR. Total routes are $\lvert\mathcal{H}\rvert + P - 1 = 8$, and the
cost is wall-clock rather than passes: a deep horizon makes a later pass
dear, and on one instance a pass is stopped by the per-pass guard.

One property matters for reading the results. Refinement consumes a
post-routing layout, and routing destroys the per-core free-slot
reservation $\ell_0$ was built with (Section~\ref{sec:sfc-layout}), so an
intermediate pass may start from a layout worse than $\ell_0$ for
\PORTER{}'s capacity machinery. Individual passes therefore regress on
some instances; only the best-pass rule keeps that from surfacing.

\subsection{Comparison Protocols and Budgets}
\label{sec:budgets}

Table~\ref{tab:budgets} states, for every configuration that appears in
the results, where its initial layouts come from, how many routing passes
it spends, and what it selects on. We use the names in its first column
throughout.

\begin{table*}[t]
\centering
\caption{Every configuration in Section~\ref{sec:results}: layout source,
  layout count $\lvert\mathcal{L}\rvert$, passes per layout $P$, routes, and
  what is minimised. Routes are $\lvert\mathcal{L}\rvert\!\cdot\!P$, plus
  $\lvert\mathcal{H}\rvert{-}1$ probes per seed where a horizon is selected.
  \appLayoutSearch{} gives the one configuration not listed.}
\label{tab:budgets}
\setlength{\tabcolsep}{6pt}
\small
\begin{tabular}{@{}l l r r r l@{}}
\toprule
Name & Layout source & $\lvert\mathcal{L}\rvert$ & $P$ &
routes & selects on \\
\midrule
\multicolumn{6}{@{}l}{\emph{Own protocol} (RQ1, RQ4)} \\
\TeleSABRE{}      & Hungarian, seeds 0--2 & 3 & 1 & 3 & best EPR \\
\dSABRE{}         & SabreLayout, seeds 0--2 & 3 & 3 & 9 & best EPR \\
\PORTER{}         & temporal locality-greedy & 1 & 4 & 4 & best EPR \\
\PORTER{}+R       & temporal locality-greedy & 1 & 5 & 8 & best forward pass \\
\midrule
\multicolumn{6}{@{}l}{\emph{Shared layout} (RQ2) --- both compilers routed from one external layout} \\
\dSABRE{} shared  & \TeleSABRE{}'s or pytket-dqc's & 1 & 3 & 3 & best pass \\
\PORTER{} shared  & the identical layout & 1 & 4 & 4 & best EPR \\
\PORTER{}+R shared& the identical layout & 1 & 5 & 8 & best forward pass \\
\midrule
\multicolumn{6}{@{}l}{\emph{Cross layout} (RQ2) --- each compiler on the other's own layout} \\
\dSABRE{} on \PORTER{} & \PORTER{}'s layout & 1 & 3 & 3 & best pass \\
\PORTER{} on \dSABRE{} & \dSABRE{}'s candidates & 3 & 4 & 12 & best EPR \\
\midrule
\multicolumn{6}{@{}l}{\emph{Ablation} (RQ3) --- one fixed layout, one fixed horizon} \\
\PORTER{} ablation & temporal locality-greedy & 1 & 1 & 1 & --- \\
\bottomrule
\end{tabular}
\end{table*}

Two notions of a matched budget are in play and neither is universally
fair, so we report against both. Fixing the \emph{initial layout} routes
both compilers from one identical mapping --- the shared-layout rows ---
which controls the layout but not the effort spent on it: \PORTER{} spends
four routes on its candidate horizons to \dSABRE{}'s three passes. Equal
\emph{candidate-layout sets} instead hand \PORTER{} the very layouts
\dSABRE{} generated for itself, so both choose their best from the same
three starting points; that equalises the search over layouts rather than
the routing spent on each. Wall-clock is not equalised anywhere, but every
table reports complete compile time so a reader can renormalise. \PORTER{}
both spends more routes in the shared-layout rows and costs more per route
(Section~\ref{sec:frontier}), so neither comparison understates it.

\subsection{Metrics and Aggregation}
\label{sec:metrics}

The primary metric is EPR pairs consumed: \TeleSABRE{} counts teledata
plus telegate operations, \dSABRE{} counts teledata, and \PORTER{} counts
teleports plus telegates plus telegate packets, one pair per packet, all
three the same physical resource. Intra-core SWAP counts and compile time
are secondary and are never traded against EPR by any selection rule.
Aggregates are unweighted geometric means over circuits, so each circuit
counts equally; percentage changes are $100(b/a-1)$ against the stated
baseline and ratios are $b/a$, both computed from those means rather than
averaged per circuit. Failures are never imputed: a configuration that
does not complete an instance is excluded from every mean it appears in,
along with that instance's other columns where the comparison is paired,
so a mean labelled $(n)$ is over exactly the $n$ circuits all its columns
completed --- which flatters a baseline that fails its hardest instances,
and is why completion counts accompany every such mean. Compile times are
wall-clock on one otherwise idle machine (Apple M2, macOS~26.5,
CPython~3.13), measured once per configuration with no warm-up, so they
are indicative and should be read as ratios between protocols rather than
as absolute costs.
